\documentclass[10pt]{article}
\usepackage{tmlr}

\usepackage{amsmath,amssymb,amsthm}
\usepackage{booktabs}
\usepackage{array}
\usepackage{siunitx}
\usepackage{graphicx}
\usepackage{caption}
\usepackage{microtype}
\usepackage{hyperref}
\usepackage{url}

% Figures are drawn with TikZ/PGFPlots from the exported data tables, so they
% are true vector art set in the document's own font rather than imported
% raster or foreign-font artwork.
\usepackage{tikz}
\usepackage{pgfplots}
\pgfplotsset{compat=1.18}
\usepgfplotslibrary{colormaps}

% Caption numbers, axis limits and colour limits emitted by
% scripts/tmlr_figdata.py alongside the figure data they describe.
\newcommand{\atlasN}{50{,}558}
\newcommand{\atlasBranches}{84}
\newcommand{\atlasGridRe}{180}
\newcommand{\atlasGridIm}{136}
\newcommand{\atlasOccupied}{20}
\newcommand{\atlasMetaMax}{2.1614}

\newcommand{\psAxmin}{1.425000}
\newcommand{\psAxmax}{2.025000}
\newcommand{\psAymin}{-0.625000}
\newcommand{\psAymax}{-0.025000}
\newcommand{\psAzmin}{-5.1802}
\newcommand{\psAzmax}{-2.7319}
\newcommand{\psGrid}{61}
\newcommand{\psBxmin}{2.459413}
\newcommand{\psBxmax}{3.059413}
\newcommand{\psBymin}{-1.723038}
\newcommand{\psBymax}{-1.123038}
\newcommand{\psBzmin}{-7.2771}
\newcommand{\psBzmax}{-4.9449}
\newcommand{\psGap}{0.5116}
\newcommand{\psPartnerIm}{-2.055}


% ---------------------------------------------------------------- tables --
% One look for every table: rules from booktabs only, no vertical lines,
% numerals decimal-aligned, a little air between rows.
\sisetup{detect-all, table-number-alignment = center, group-digits = none}
\renewcommand{\arraystretch}{1.18}
\captionsetup[table]{skip=6pt}
\captionsetup{font=small, labelfont=bf}
\newcommand{\thead}[1]{\multicolumn{1}{c}{#1}}

\theoremstyle{plain}
\newtheorem{proposition}{Proposition}
\newtheorem{lemma}{Lemma}
\newtheorem{theorem}{Theorem}

\DeclareMathOperator{\spec}{spec}
\DeclareMathOperator{\tr}{tr}
\newcommand{\Rey}{\operatorname{Re}}
\newcommand{\Imy}{\operatorname{Im}}
\newcommand{\dmin}{d^{(\mathrm{sample})}_{\min}}

% The trailing star is the marker for the author-disclosure footnote set
% immediately after \maketitle; it is suppressed in the PDF bookmark string.
\title{Azimuthal-Sector Protection and Exceptional-Point Search\\
in Five-Dimensional Myers--Perry Quasinormal
Spectra\texorpdfstring{\textsuperscript{*}}{}}

\author{\name Anonymous Authors \email anonymous@example.org \\
        \addr Anonymous Institution}

\def\month{08}
\def\year{2026}
\def\openreview{\url{https://openreview.net/forum?id=XXXX}}

% Document properties for the reader's PDF viewer.  The author field stays
% empty: this is a double-blind submission, and the title is not identifying.
\hypersetup{
  pdftitle={Azimuthal-Sector Protection and Exceptional-Point Search in
            Five-Dimensional Myers-Perry Quasinormal Spectra},
  pdfauthor={},
  pdfsubject={Submission to Transactions on Machine Learning Research},
  hidelinks,
}

\begin{document}
\maketitle

% Author disclosure carried over from the submitted version, kept as the
% starred title footnote it was originally set as.
\renewcommand{\thefootnote}{\fnsymbol{footnote}}
\footnotetext[1]{OpenAI ChatGPT (GPT-5.6 Sol) was used as an assistive tool for
editorial organization, literature-search planning, code prototyping,
consistency checking, and figure-layout refinement. The authors independently
reviewed and approved all scientific claims, mathematical arguments, numerical
results, references, figures, and text and take full responsibility for the
submitted work.}
\renewcommand{\thefootnote}{\arabic{footnote}}
\setcounter{footnote}{0}

\begin{abstract}
We study whether the enhanced symmetry of equal-magnitude-spin five-dimensional
Myers--Perry black holes can organize defective degeneracies in the
massive-scalar quasinormal spectrum. Away from the mass-threshold branch
points, we assume a local holomorphic Fredholm realization of the quasinormal
pencil that commutes with the surviving azimuthal $T^2$ action. The associated
Fourier projectors reduce the pencil into invariant $(m_1,m_2)$ sectors and
project Jordan chains sector by sector. We prove that a coincidence of
frequencies from distinct azimuthal sectors cannot increase the maximal
Jordan-chain length or the order of a resolvent pole; any defectiveness must
already occur within one fixed sector. The discrete symmetry is the order-eight
signed-permutation group generated by independent plane reflections and plane
interchange, with $a=b$ and $a=-b$ forming isometric components of the
equal-magnitude-spin locus. In diagonal sectors, analytic characteristic
invariants are even in $\delta=(a-b)/2$, so a quotient curve transverse to
$t=\delta^2=0$ lifts with quadratic tangency in the remaining coordinates. We
then search the surviving within-sector collision channel in a finite atlas of
$50{,}558$ continued mode entries across seven azimuthal sectors. No sampled or
locally refined candidate simultaneously exhibits gap closure, one-loop branch
exchange, and square-root splitting. The smallest stored branch gap is
$0.5116$. The numerical result is therefore a finite-atlas non-detection, not a
continuum exclusion theorem.
\end{abstract}

\section{Introduction}

Quasinormal modes (QNMs) are resonances of black-hole perturbation operators
and poles of the retarded Green function
\citep{kokkotas1999,berti2009,konoplya2011,berti2026}. Their spectral problem
is non-self-adjoint because radiation leaves through the horizon and infinity.
Consequently, eigenvalues may coalesce together with their eigenvectors at
exceptional points (EPs) \citep{kato1995,bergholtz2021,ashida2020}.

Let a perturbation problem be represented locally by an operator pencil
$T(\omega)$. If an isolated eigenvalue $\omega_*$ has maximal Jordan-chain
length $q$, the Keldysh expansion gives a pole of order $q$ in
$T(\omega)^{-1}$. Its inverse transform contributes
\begin{equation}
\psi(t)\sim\bigl(A_0+A_1t+\cdots+A_{q-1}t^{\,q-1}\bigr)e^{-i\omega_*t}.
\label{eq:keldysh}
\end{equation}
At an EP2 the generalized mode therefore contains a $te^{-i\omega_*t}$ term.
Ordinary eigenvectors alone are incomplete at the EP, while a Jordan chain
supplies the appropriate generalized spectral data. Near a generic EP2,
eigenvalue splitting is proportional to the square root of a transversal
perturbation \citep{kato1995,macedo2026}.

Exceptional and near-exceptional structures have been identified in several
black-hole spectra, including massive-scalar Kerr and model ringdown systems
\citep{cavalcante2024,cavalcante2026,motohashi2025,macedo2026,cao2025}. Recent
black-hole resonance analyses also emphasize the local Riemann-surface
description of modes near exceptional points \citep{yang2025}. The
five-dimensional Myers--Perry geometry is attractive for three reasons. The
massive Klein--Gordon equation separates for two arbitrary spins
\citep{frolov2003,murata2008,huang2025}; the family has two rotational
parameters plus the scalar mass; and the continuous isometry is enhanced to
$\mathbb{R}\times U(2)$ when the two angular momenta have equal magnitude
\citep{myers1986,myers2012}.

Our analysis has two complementary parts. First, we isolate the exact symmetry
mechanism that constrains defectiveness. Second, we search numerically for the
within-sector collisions that the theorem leaves open. The numerical claim is
intentionally restricted to the computed atlas: without a continuum enclosure,
a certified interpolation bound, or a rigorous radial-truncation estimate, a
sampled minimum cannot be promoted to a global lower bound over a continuous
parameter box.

The main results are as follows.
\begin{enumerate}
\item Under a strongly continuous torus action commuting with a locally
  holomorphic Fredholm pencil, cross-sector frequency coincidences cannot
  create longer Jordan chains.
\item The exact discrete symmetry is the order-eight group generated by
  independent plane reflections and plane interchange. The loci $a=b$ and
  $a=-b$ are isometric.
\item In sectors fixed by plane interchange, analytic characteristic invariants
  are even in $\delta=(a-b)/2$. A transversal EP curve in $t=\delta^2$ lifts to
  a smooth parity-symmetric physical curve whose non-$\delta$ coordinates vary
  quadratically through $\delta=0$.
\item In the numerical atlas no candidate displays coalescence, branch
  exchange, and square-root splitting simultaneously. This is a finite-atlas
  non-detection result, not a continuum theorem.
\end{enumerate}

\section{Geometry and the separated scalar problem}

\subsection{Metric, horizons, and thermodynamics}

In Boyer--Lindquist-type coordinates $(t,r,\theta,\phi,\psi)$, with
$0\leq\theta\leq\pi/2$ and $\phi,\psi$ $2\pi$-periodic, the five-dimensional
Myers--Perry metric is
\begin{align}
ds^2 &= -dt^2+\frac{M}{\Sigma}
  \bigl(dt-a\sin^2\theta\,d\phi-b\cos^2\theta\,d\psi\bigr)^2
  +\frac{\Sigma}{\Delta}dr^2+\Sigma\,d\theta^2 \nonumber\\
 &\quad +(r^2+a^2)\sin^2\theta\,d\phi^2+(r^2+b^2)\cos^2\theta\,d\psi^2 ,
\label{eq:metric}
\end{align}
where
\begin{equation}
\Sigma=r^2+a^2\cos^2\theta+b^2\sin^2\theta,\qquad
\Pi=(r^2+a^2)(r^2+b^2),\qquad
\Delta=\frac{\Pi-Mr^2}{r^2}.
\end{equation}
The mass parameter $M$ has dimensions of length squared. We set $M=1$ in
numerical expressions. The ADM quantities are
\begin{equation}
M_{\mathrm{ADM}}=\frac{3\pi M}{8},\qquad
J_\phi=\tfrac{2}{3}M_{\mathrm{ADM}}a,\qquad
J_\psi=\tfrac{2}{3}M_{\mathrm{ADM}}b.
\end{equation}
With $z=r^2$, the horizon polynomial is
\begin{equation}
P(z)=(z+a^2)(z+b^2)-Mz=(z-z_+)(z-z_-),
\end{equation}
with
\begin{equation}
z_\pm=\tfrac{1}{2}\Bigl[M-a^2-b^2\pm\sqrt{(M-a^2-b^2)^2-4a^2b^2}\Bigr].
\end{equation}
A horizon exists when $M\geq(|a|+|b|)^2$. We use
\begin{equation}
\mathcal{E}=M-(|a|+|b|)^2
\end{equation}
as an extremality measure. It is also useful to introduce
\begin{equation}
s=\frac{a+b}{2},\qquad \delta=\frac{a-b}{2}.
\end{equation}
The equal-magnitude-spin locus $a^2=b^2$ has two coordinate components,
$\delta=0$ and $s=0$. A reflection of either rotation plane maps one component
to the other, so they are isometric descriptions of the same equal-magnitude
condition.

The horizon angular velocities, surface gravity, temperature, area, and entropy
are
\begin{gather}
\Omega_a=\frac{a}{z_++a^2},\qquad \Omega_b=\frac{b}{z_++b^2},\\
\kappa=\frac{r_+(z_+-z_-)}{(z_++a^2)(z_++b^2)},\qquad T_H=\frac{\kappa}{2\pi},\\
A_H=\frac{2\pi^2}{r_+}(z_++a^2)(z_++b^2),\qquad S_{\mathrm{BH}}=\frac{A_H}{4}.
\end{gather}
Representative horizon data used to check the parameterization are collected in
Table~\ref{tab:horizons}.

\begin{table}[!ht]
\centering
\caption{Representative horizon data in units $M=1$. The final two rows are
related by interchange of the rotation planes, which swaps $\Omega_a$ and
$\Omega_b$.}
\label{tab:horizons}
\begin{tabular}{
  S[table-format=1.2]
  S[table-format=+1.2]
  S[table-format=1.4]
  S[table-format=1.5]
  S[table-format=1.4]
  S[table-format=1.4]
}
\toprule
\thead{$s$} & \thead{$\delta$} & \thead{$r_+$} & \thead{$T_H$} &
\thead{$\Omega_a$} & \thead{$\Omega_b$} \\
\midrule
  0.00 & +0.00 & 1.0000 & 0.15915 & 0.0000 & 0.0000 \\
  0.20 & +0.00 & 0.9583 & 0.15222 & 0.2087 & 0.2087 \\
  0.42 & +0.00 & 0.7713 & 0.11196 & 0.5445 & 0.5445 \\
  0.34 & +0.20 & 0.8249 & 0.12966 & 0.5556 & 0.2000 \\
  0.42 & +0.20 & 0.7296 & 0.10849 & 0.6764 & 0.3789 \\
  0.42 & -0.20 & 0.7296 & 0.10849 & 0.3789 & 0.6764 \\
\bottomrule
\end{tabular}

\end{table}

\subsection{Separation}

For a scalar of mass $\mu$, $(\Box-\mu^2)\Phi=0$, take
\begin{equation}
\Phi=e^{-i\omega t+im_1\phi+im_2\psi}R(r)S(\theta),\qquad m_1,m_2\in\mathbb{Z}.
\label{eq:ansatz}
\end{equation}
The separated equations are
\begin{align}
0&=\frac{1}{\sin\theta\cos\theta}\frac{d}{d\theta}
   \Bigl(\sin\theta\cos\theta\frac{dS}{d\theta}\Bigr)
   +\Bigl[\Lambda+(\omega^2-\mu^2)\bigl(a^2\cos^2\theta+b^2\sin^2\theta\bigr)
   -\frac{m_1^2}{\sin^2\theta}-\frac{m_2^2}{\cos^2\theta}\Bigr]S,
\label{eq:angular}\\
0&=\frac{1}{r}\frac{d}{dr}\Bigl(r\Delta\frac{dR}{dr}\Bigr)
   +\Bigl[\frac{W^2}{r^4\Delta}-\frac{G^2}{r^2}-(a^2+b^2)\omega^2
   +2\omega(am_1+bm_2)-\mu^2r^2-\Lambda\Bigr]R,
\label{eq:radial}
\end{align}
where
\begin{equation}
W=\Pi\omega-m_1a(r^2+b^2)-m_2b(r^2+a^2),\qquad G=ab\omega-am_2-bm_1.
\end{equation}
At $a=b=0$, $\Lambda=\ell(\ell+2)$. These equations agree algebraically with the
arbitrary-two-spin formulation in \citet{huang2025}, after accounting for the
convention used for the separation constant.

Set
\begin{equation}
\widehat{A}=\Lambda+(\omega^2-\mu^2)b^2,\qquad
c^2=(\omega^2-\mu^2)(a^2-b^2)=4s\delta(\omega^2-\mu^2).
\end{equation}
When $a^2=b^2$, $c^2=0$ and the angular operator is the round scalar Laplacian
on $S^3$,
\begin{equation}
\widehat{A}=\ell(\ell+2),\qquad \ell=2n+|m_1|+|m_2|,\qquad n\in\mathbb{Z}_{\geq0}.
\end{equation}
The angular eigenspace has dimension $(\ell+1)^2$. For the full rotating
spacetime, the relevant continuous isometry on either equal-magnitude component
is $U(2)$ rather than the full static $SO(4)$ symmetry \citep{myers2012}.

At the outer horizon,
\begin{equation}
\sigma_+=\frac{\omega-m_1\Omega_a-m_2\Omega_b}{2\kappa}=\frac{\omega-\omega_c}{2\kappa},
\qquad \omega_c=m_1\Omega_a+m_2\Omega_b,
\end{equation}
and the QNM boundary condition selects
\begin{equation}
R\sim(z-z_+)^{-i\sigma_+}.
\label{eq:horizonbc}
\end{equation}
At large $r$,
\begin{equation}
R''+\frac{3}{r}R'+\bigl(\Omega^2+O(r^{-2})\bigr)R=0,\qquad
\Omega=\sqrt{\omega^2-\mu^2},
\end{equation}
so the outgoing branch is
\begin{equation}
R\sim r^{-3/2}e^{+i\Omega r}\bigl(1+O(r^{-1})\bigr).
\label{eq:infinitybc}
\end{equation}
There is no $1/r$ Coulomb tail in this five-dimensional radial equation.

\subsection{Superradiance and the local analytic pencil}

For a real-frequency scattering state, propagation at infinity requires
$\Omega\in\mathbb{R}$, hence $\omega>\mu$ for positive frequency. The horizon
flux has the sign of $\omega-\omega_c$. Scattering amplification therefore
occurs for
\begin{equation}
\mu<\omega<\omega_c.
\end{equation}
For a quasibound state, the field instead decays at infinity and a superradiant
instability may occur when $0<\Rey\omega<\omega_c$ together with
$\Rey\omega<\mu$ \citep{detweiler1980,dolan2007}. The ingoing horizon condition
\eqref{eq:horizonbc} remains the physical horizon condition on both sides of
$\omega=\omega_c$; what changes is the flux sign, not the definition of ingoing
propagation.

The boundary condition at infinity contains the square root
$\Omega=\sqrt{\omega^2-\mu^2}$. Accordingly, a holomorphic pencil description is
local: one chooses a sheet of the square root and works away from the threshold
branch points $\omega=\pm\mu$. On such a neighborhood the coupled
angular-radial problem may be written
\begin{equation}
T(\omega;s,\delta,\mu,m_1,m_2)\,u=0,
\end{equation}
with $\Lambda(\omega)$ recomputed at every evaluation. Derivatives with respect
to $\omega$ must be total derivatives; freezing $\Lambda$ drops the term
$(\partial_\Lambda F)(d\Lambda/d\omega)$ and can create spurious multiple-root
diagnostics.

\section{Azimuthal-sector protection theorem}

\subsection{Torus action and projectors}

Neither the metric nor the boundary conditions depend explicitly on $\phi$ or
$\psi$. Define the torus action
\begin{equation}
(R_{\alpha\beta}u)(t,r,\theta,\phi,\psi)=u(t,r,\theta,\phi-\alpha,\psi-\beta).
\label{eq:torusaction}
\end{equation}
A mode proportional to $e^{im_1\phi+im_2\psi}$ transforms with character
$e^{-i(m_1\alpha+m_2\beta)}$. Therefore the projector onto the sector labeled by
the ansatz \eqref{eq:ansatz} is
\begin{equation}
\Pi_{m_1m_2}=\frac{1}{(2\pi)^2}\int_0^{2\pi}\!\!\int_0^{2\pi}
  e^{+i(m_1\alpha+m_2\beta)}R_{\alpha\beta}\,d\alpha\,d\beta.
\label{eq:projector}
\end{equation}
The plus sign in Eq.~\eqref{eq:projector} is essential for the convention
\eqref{eq:torusaction}; the opposite sign projects onto $(-m_1,-m_2)$.

\begin{proposition}[Invariant sector decomposition]
Let $X$ be a Banach space on which $R_{\alpha\beta}$ is a uniformly bounded,
strongly continuous representation of $T^2$. Suppose that $T(\omega)$ and its
local $\omega$ derivatives commute with the action and that the operator domain
is invariant. Then the projectors \eqref{eq:projector} are bounded, pairwise
annihilating idempotents that commute with the pencil. Moreover,
\begin{equation}
\Pi_{m_1m_2}u=0 \ \text{ for all } (m_1,m_2)
\quad\Longrightarrow\quad u=0.
\end{equation}
\end{proposition}

\begin{proof}
The Bochner integral exists by uniform boundedness. Idempotence and pairwise
annihilation follow from character orthogonality. To show separation of points
without assuming unconditional Fourier-series convergence, convolve the orbit
map with the two-dimensional Fej\'er kernel. The Fej\'er sums are finite
positive combinations of the projectors and converge in norm to $u$ by strong
continuity. If all projected components vanish, every Fej\'er sum vanishes and
hence $u=0$.
\end{proof}

The statement is local in spectral parameter. We assume a weighted or
hyperboloidal realization for which the quasinormal pencil is holomorphic and
Fredholm on a chosen sheet away from $\omega=\pm\mu$, and for which the torus
action preserves the operator domain. For an isolated characteristic value, the
root subspace is finite dimensional; its torus representation therefore
contains only finitely many characters \citep{guttel2017}. No global Fredholm
construction at the mass threshold is asserted.

\subsection{Jordan chains and pole order}

A Jordan chain $(v_0,\dots,v_{q-1})$, with $v_0\neq0$, satisfies
\begin{equation}
\sum_{i=0}^{j}\frac{1}{i!}T^{(i)}(\omega_*)v_{j-i}=0,\qquad j=0,\dots,q-1.
\label{eq:chain}
\end{equation}

\begin{lemma}
Under the preceding hypotheses, every Jordan chain projects componentwise to
Jordan chains of the sector restrictions. At least one projected chain has the
same length, and a chain in one sector extends by zero to the full space. Thus
the maximal chain length of $T$ is the maximum of the sectorwise maximal chain
lengths.
\end{lemma}

\begin{proof}
Apply $\Pi_{m_1m_2}$ to every equation in \eqref{eq:chain}. Commutation with
$T^{(i)}$ gives the sector chain equations. Since $v_0\neq0$, the separation
property implies that at least one $\Pi_{m_1m_2}v_0$ is nonzero. The converse
extension is immediate.
\end{proof}

\begin{theorem}[Cross-sector no-go]
Assume the local holomorphic Fredholm realization described above, and let
$\omega_*$ be an isolated characteristic value. The maximal Jordan-chain length
at $\omega_*$ equals the maximum of the maximal chain lengths of the sector
restrictions. Hence a frequency coincidence between distinct azimuthal sectors
cannot, by itself, create defectiveness, raise the resolvent-pole order, or
generate an additional secular power of $t$ in Eq.~\eqref{eq:keldysh}. Any
chain of length greater than one is contained in a single fixed $(m_1,m_2)$
sector.
\end{theorem}

The wording does not exclude a frequency shared by two sectors when one sector
independently contains an EP. In that case the defect belongs to the single
sector and is not caused by the cross-sector coincidence.

\subsection{Equal-magnitude-spin multiplets}

On $a=b$ the radial equation depends on the azimuthal labels through
$m=m_1+m_2$. At fixed $(\ell,m)$, the degenerate labels are
\begin{equation}
\mathcal{M}_{\ell,m}=\bigl\{(m_1,m_2):m_1+m_2=m,\ |m_1|+|m_2|\leq\ell,\
\ell-|m_1|-|m_2|\in2\mathbb{Z}_{\geq0}\bigr\}.
\end{equation}
For each admissible $m=-\ell,-\ell+2,\dots,\ell$, one finds
$|\mathcal{M}_{\ell,m}|=\ell+1$. Summing over the $\ell+1$ values of $m$ gives
the angular degeneracy $(\ell+1)^2$. Members of a fixed $\mathcal{M}_{\ell,m}$
occupy distinct azimuthal sectors, so their coincidence is semisimple across
sectors unless a participating sector is independently defective.

The continuation confirms this directly. On the slice $\ell=2$, $N=0$,
$s=0.20$, $\mu=0.6$, the $(1,1)$ and $(2,0)$ branches agree at $\delta=0$ to
$\Rey\omega=1.6835979353$ in both sectors, coinciding to $4\times10^{-15}$,
and separate to $1.6892638$ and $1.8006703$ at $\delta=0.20$. The diagonal
sector is even in $\delta$ to machine precision, returning $1.6892638$ at
$\delta=-0.20$ as well, while the off-diagonal sector maps to its partner
$(0,2)$ and returns $1.5994653$. No spectral surface is inferred outside this
computed slice. The theorem protects against cross-sector hybridization, not
against spectral splitting.

The component $a=-b$ is obtained from $a=b$ by reflecting one azimuthal plane.
Its radial multiplets are equivalently organized by $m_1-m_2$ after the
corresponding label reflection.

\section{Discrete symmetries and local exceptional topology}

\subsection{The full signed-permutation group}

In addition to interchange of the two rotation planes, the scalar problem is
invariant under independent orientation reversals:
\begin{align}
A&:\ (a,m_1,\phi)\mapsto(-a,-m_1,-\phi),\\
B&:\ (b,m_2,\psi)\mapsto(-b,-m_2,-\psi),\\
E&:\ (a,b,m_1,m_2,\theta,\phi,\psi)\mapsto
   \bigl(b,a,m_2,m_1,\tfrac{\pi}{2}-\theta,\psi,\phi\bigr).
\end{align}
They generate
\begin{equation}
G=\langle A,B,E\rangle\simeq(\mathbb{Z}_2\times\mathbb{Z}_2)\rtimes S_2,
\qquad |G|=8.
\end{equation}
The subgroup $\{1,E,AB,EAB\}$ is a Klein four-group, but it is not the full
discrete symmetry. In $(s,\delta)$ coordinates,
\begin{align}
E&:(s,\delta)\mapsto(s,-\delta), & AB&:(s,\delta)\mapsto(-s,-\delta),\\
A&:(s,\delta)\mapsto(-\delta,-s), & B&:(s,\delta)\mapsto(\delta,s).
\end{align}
Table~\ref{tab:stabilizers} lists the nontrivial fixed-sector symmetries. The
equalities are naturally setwise statements for the spectrum; an isolated simple
branch can be labeled to satisfy them pointwise.

\begin{table}[t]
\centering
\caption{Fixed-sector stabilizers of the signed-permutation symmetry. Spectral
equalities are setwise; an isolated simple branch can be labeled pointwise.}
\label{tab:stabilizers}
\begin{tabular}{cccl}
\toprule
\thead{Sector} & \thead{Stabilizer} & \thead{Parameter action} &
\thead{Spectral relation} \\
\midrule
$m_1=m_2$  & $E$    & $(s,\delta)\mapsto(s,-\delta)$
           & $\spec_{m,m}(s,\delta)=\spec_{m,m}(s,-\delta)$ \\
$m_1=-m_2$ & $EAB$  & $(s,\delta)\mapsto(-s,\delta)$
           & $\spec_{m,-m}(s,\delta)=\spec_{m,-m}(-s,\delta)$ \\
$m_1=0$    & $A$    & $(s,\delta)\mapsto(-\delta,-s)$
           & $\spec_{0,m_2}(s,\delta)=\spec_{0,m_2}(-\delta,-s)$ \\
$m_2=0$    & $B$    & $(s,\delta)\mapsto(\delta,s)$
           & $\spec_{m_1,0}(s,\delta)=\spec_{m_1,0}(\delta,s)$ \\
$(0,0)$    & $G$    & all signed permutations
           & full induced invariance \\
\bottomrule
\end{tabular}
\end{table}

\subsection{Discriminant, semisimple coincidences, and the
\texorpdfstring{$t=\delta^2$}{t = delta squared} quotient}

A local two-mode spectral subspace may be represented by a $2\times2$ matrix
$H(p)$. The similarity-invariant quantities are
\begin{equation}
\mathcal{T}=\tr H,\qquad \mathcal{D}=(\tr H)^2-4\det H,
\end{equation}
with eigenvalues
\begin{equation}
\omega_\pm=\tfrac{1}{2}\bigl(\mathcal{T}\pm\sqrt{\mathcal{D}}\bigr).
\end{equation}
The equation $\mathcal{D}=0$ gives a repeated algebraic eigenvalue. It is not,
by itself, equivalent to an EP: the special case $H=\omega_*I$ is a semisimple
double eigenvalue. A simple zero of $\mathcal{D}$ in a transversal direction is
generically defective and yields the square-root Puiseux behavior
\begin{equation}
\omega_\pm(\tau)=\omega_*+c_1\tau\pm c_{1/2}\tau^{1/2}+O(\tau^{3/2}),
\qquad c_{1/2}\neq0.
\end{equation}
Geometric multiplicity or a Jordan-chain test must still be checked at the
repeated root.

In a diagonal sector, the characteristic polynomial and its analytic invariants
are even in $\delta$. Individual eigenvalue labels need not be analytic or
separately even at an EP. Let $t=\delta^2$ and suppose an EP component in
quotient coordinates $(s,t,\mu)$ meets $t=0$ transversely. With a local curve
parameter $\tau$,
\begin{equation}
t=c\tau+O(\tau^2),\qquad c\neq0.
\end{equation}
For any other smooth coordinate $q$ along the curve (for example $s$ or $\mu$),
write $q=q_*+A\tau+O(\tau^2)$. Eliminating $\tau$ gives
\begin{equation}
\tau=\frac{\delta^2}{c}+O(\delta^4),\qquad
q-q_*=\frac{A}{c}\delta^2+O(\delta^4).
\label{eq:lift}
\end{equation}
The two quotient preimages therefore join into a single smooth
parity-symmetric curve in the full physical parameter space. At $\delta=0$ the
tangent direction is along the $\delta$ coordinate, while the remaining
coordinates have a parabolic profile. The quotient locus has an endpoint at
$t=0$; the physical curve does not terminate there. Equation~\eqref{eq:lift} is
the local normal form of this lift.

\section{Exceptional-point diagnostics}

A credible negative search requires diagnostics that are invariant under
harmless reformulations of the spectral condition and that are calibrated on
positive and negative controls.

\subsection{Admissible signatures}

For a generic EP2, the following signatures are complementary:
\begin{enumerate}
\item independently tracked eigenvalue separation tends to zero under local
  refinement;
\item one circuit around the candidate transposes the two branches, and two
  circuits restore them;
\item the local splitting is better fit by $|\tau|^{1/2}$ than by $|\tau|$ or a
  nonzero constant \citep{nennig2020};
\item the repeated root has geometric multiplicity one, or an associated vector
  solves the Jordan-chain equation;
\item the result is stable under discretization, precision, and a numerically
  independent formulation.
\end{enumerate}
For a meromorphic scalar condition $F(\omega)$, the argument principle gives
\begin{equation}
\frac{1}{2\pi i}\oint_C\frac{F'(\omega)}{F(\omega)}\,d\omega=N_Z-N_P .
\end{equation}
Zeros and poles must therefore be separated, rather than interpreting
$N_Z-N_P$ as a zero count when poles are present.

A matrix-based nonlinear-eigenvalue condition indicator may be written
\begin{equation}
\kappa_T(\omega_*)=\frac{\|u\|\,\|v\|}{|\langle v,T'(\omega_*)u\rangle|},
\label{eq:kappa}
\end{equation}
where $u$ and $v$ are right and left null vectors and $v$ belongs to the dual
space in the operator setting. In a fixed matrix normalization, growth of
$\kappa_T$ is a useful qualitative warning. It is not an absolute
distance-to-EP bound: under $T\mapsto g(\omega)T$ with $g(\omega_*)\neq0$,
Eq.~\eqref{eq:kappa} scales by $|g(\omega_*)|^{-1}$ unless a normalization
convention is fixed.

\subsection{Normalization-sensitive diagnostics}

The derivative $|F'(\omega_*)|$ of an arbitrarily normalized scalar spectral
condition has no invariant magnitude. Likewise, if
\begin{equation}
F=a_1(\omega-\omega_1)+a_2(\omega-\omega_1)^2+\cdots,
\end{equation}
then under $F\mapsto gF$, with $g=g_0+g_1(\omega-\omega_1)+\cdots$,
\begin{equation}
\frac{\widetilde{a}_1}{\widetilde{a}_2}
=\frac{a_1/a_2}{1+(g_1/g_0)(a_1/a_2)} .
\end{equation}
The coefficient ratio is only asymptotically insensitive to normalization. More
importantly, for a meromorphic continued fraction it responds to the nearest
zero or pole. It cannot serve as a rigorous spectral separation bound.

Table~\ref{tab:calibration} records the calibration of these diagnostics on
exactly solvable controls, computed with the same routines that are applied to
the atlas. The fitted Puiseux exponent cleanly separates the three local laws,
and the coefficient-ratio diagnostic inflates by an order of magnitude under a
rescaling that leaves the zero count exactly invariant. Calibration and
candidate evidence play different roles in an EP classification and are kept
apart throughout.

\begin{table}[t]
\centering
\caption{Diagnostic calibration on exactly solvable controls. The fitted
exponent separates square-root, linear, and floored splitting; the
normalization test shows that the coefficient ratio is not an invariant
proximity measure, whereas the argument-principle zero count is.}
\label{tab:calibration}
\begin{tabular*}{0.88\linewidth}{@{\extracolsep{\fill}}l c l@{}}
\toprule
\multicolumn{3}{@{}l}{\emph{Local splitting laws:} fitted Puiseux exponent $p$ over $t\in[10^{-8},10^{-4}]$}\\
\midrule
\thead{Model} & \thead{$p$} & \thead{Classification}\\
\midrule
  square root, $\omega_\pm=\pm\sqrt{t}$ & 0.500 & EP2 \\
  linear, $\omega_\pm=\pm 3t$ & 1.000 & ordinary crossing \\
  floored, $2\sqrt{t^2+g^2}$ & 0.000 & avoided crossing \\
\bottomrule
\end{tabular*}
\\[1.4ex]
\begin{tabular*}{0.88\linewidth}{@{\extracolsep{\fill}}l c c c@{}}
\toprule
\multicolumn{4}{@{}l}{\emph{Normalization sensitivity} under $F\mapsto gF$, $g=10^{7}e^{3\omega}$}\\
\midrule
\thead{$d$} & \thead{measured distortion} & \thead{$|1+(g'/g)d|^{-1}$} & \thead{zero count}\\
\midrule
  $1\times10^{-4}$ & 1.000 & 1.000 & 2 \\
  $1\times10^{-2}$ & 1.031 & 1.031 & 2 \\
  $3\times10^{-1}$ & 10.000 & 10.000 & 2 \\
\bottomrule
\end{tabular*}

\end{table}

\section{Numerical formulations and validation}

The numerical calculation combines recurrence-based and recurrence-free
formulations. Implementation details, atlas census information, and diagnostic
checks are summarized in the Supplementary Material.

\subsection{Continued fraction}

Introduce the compact radial coordinate $x=(z-z_+)/(z-z_-)$ and factor the
ingoing outer-horizon behavior and outgoing large-$r$ behavior from the radial
solution. The remaining analytic factor is expanded in powers of $x$. The
radial Frobenius expansion yields a finite-band recurrence
\begin{equation}
\sum_{j=0}^{p}c_j(n)\,a_{n-j}=0,\qquad p+1=13,
\end{equation}
for generic two-spin parameters, reducing to width $9$ in the static case.
Gaussian elimination gives a three-term recurrence and a Leaver continued
fraction \citep{leaver1985}. The recurrence coefficients are generated
algorithmically by Fourier coefficient extraction as described in the
Supplementary Material. Because the coefficient extraction and radial
truncation are not accompanied by a rigorous aliasing/truncation enclosure, the
continued-fraction calculation is used as numerical evidence rather than as a
certified spectral enclosure. The inversion index must follow the tracked
overtone: finite-depth inversions are not numerically interchangeable. The
large-$n$ ratio used for tail acceleration is
\begin{equation}
\frac{a_n}{a_{n-1}}=1+u_1n^{-1/2}+u_2n^{-1}+O(n^{-3/2}),\qquad
u_1=-\sqrt{-2c},\qquad c=i\Omega(r_+-r_-),
\end{equation}
with the square-root sign selected by minimality.

\subsection{Exterior complex scaling}

A recurrence-free formulation rotates the radial contour beyond $r_m$,
\begin{equation}
r=r_m+\varsigma e^{i\vartheta},\qquad 0\leq\varsigma\leq L.
\end{equation}
The outgoing factor decays along the contour when
\begin{equation}
\tan\vartheta>-\frac{\Imy\Omega}{\Rey\Omega},
\label{eq:scaling}
\end{equation}
for $\Rey\Omega>0$. Resonances should be stable under changes of $\vartheta$,
whereas the rotated continuum moves.

\subsection{Solver inventory and benchmarks}

The solver inventory is summarized in Table~\ref{tab:solvers}. Solvers A and B
share the same recurrence and are not independent. Solver C is the principal
recurrence-free cross-check. The validation suite reproduces the
five-dimensional Schwarzschild--Tangherlini scalar fundamental mode
\begin{equation}
\omega=0.533835574268-0.383375368513\,i \qquad (r_+=1)
\end{equation}
within $5.0\times10^{-13}$ of \citet{matyjasek2021}, exchange symmetry to
$4.2\times10^{-15}$, and three-solver agreement below $10^{-5}$ on ten
benchmark points; comparable spectral-method cross-checks are surveyed by
\citet{hatsuda2021} and \citet{stein2019}. These checks support implementation
consistency at selected points, but they do not bound the numerical error
throughout the continuous parameter domain.

\begin{table}[t]
\centering
\caption{Numerical formulations and their role in the validation strategy.
Solvers A and B share the Frobenius recurrence; C and D are recurrence-free
collocation formulations.}
\label{tab:solvers}
\begin{tabular}{c >{\raggedright\arraybackslash}p{0.265\textwidth} c
                >{\raggedright\arraybackslash}p{0.345\textwidth}}
\toprule
\thead{Solver} & \thead{Formulation} & \thead{Recurrence-free} &
\thead{Primary role / limitation} \\
\midrule
A & Gaussian-reduced Leaver continued fraction & no
  & Main continuation solver; finite-depth drift near threshold \\
B & Raw recurrence, Hill determinant, Wynn acceleration & no
  & Same recurrence family as A; not an independent cross-check \\
C & Exterior-complex-scaled Chebyshev collocation & yes
  & Independent resonance check when an admissible scaled contour exists \\
D & Multidomain near-horizon collocation & yes
  & Cross-check for ordinary modes; shares the outgoing-contour threshold
    obstruction \\
\bottomrule
\end{tabular}
\end{table}

\section{Finite-atlas search and numerical result}

\subsection{Atlas construction}

Branches are labeled by continuation history rather than instantaneous sorting.
The continuation acceptance rule for the next predictor miss is
\begin{equation}
\epsilon_{k+1}\leq\max\bigl[4\,\mathrm{median}(\epsilon_1,\dots,\epsilon_k),\
0.02\max(|\omega_{k+1}|,1)\bigr].
\end{equation}
The finite atlas contains $50{,}558$ mode entries in the seven independent
sectors
\begin{equation}
(0,0),\ (1,1),\ (-1,-1),\ (1,0),\ (2,0),\ (2,1),\ (2,2),
\end{equation}
with
\begin{equation}
\ell\in\{\ell_{\min},\ell_{\min}+2,\ell_{\min}+4\},\qquad
\ell_{\min}=|m_1|+|m_2|,\qquad N=0,1,2,3 .
\end{equation}
The declared sampled region is
\begin{equation}
0\leq s\leq0.42,\qquad |\delta|\leq0.20,\qquad 0\leq\mu\leq1.8,\qquad
r_-\leq0.2287,\qquad \mathcal{E}\geq0.294 .
\label{eq:box}
\end{equation}
Table~\ref{tab:census} summarizes the continuation census, including predictor
misses, nonconvergences, and excluded branch-label collapses.
Figure~\ref{fig:atlas} shows the stored atlas itself in the complex-frequency
plane. The four horizontal ladders are the overtones $N=0,1,2,3$; the
filamentary texture within each ladder is the trace of the continuation sweeps
in $(s,\delta,\mu)$, and the empty regions between them are genuinely empty
rather than merely unresolved.

\begin{figure}[t]
\centering
% Figure 1 -- the stored quasinormal atlas in the complex-frequency plane.
% Drawn in PGFPlots from data/atlas_density.dat: occupancy of all 50,558
% stored entries on a 180 x 136 grid, as log10(count+1).  Every entry is
% counted and nothing is smoothed, so an empty cell is exactly empty.
% The axis, ticks, labels, colour bar and annotations are typeset by PGFPlots
% in the document's own font; only the occupancy field itself is pixel data.
\begin{tikzpicture}

\pgfplotsset{
  colormap={emberdark}{
    rgb=(0.001,0.000,0.014)
    rgb=(0.232,0.060,0.438)
    rgb=(0.554,0.161,0.506)
    rgb=(0.868,0.288,0.409)
    rgb=(0.988,0.553,0.348)
    rgb=(0.987,0.991,0.750)
  },
}
% Labels sit on top of the field, so each one carries its own dark plate;
% white-on-bright is the one combination this colour map cannot survive.
\tikzset{
  plate/.style={font=\scriptsize, text=white, inner xsep=2.4pt,
                inner ysep=1.6pt, fill=black, fill opacity=0.55,
                text opacity=1, rounded corners=1pt},
  band/.style={plate, anchor=west},
  winlab/.style={plate, font=\fontsize{7}{8}\selectfont},
  win/.style={draw=white, line width=0.6pt, rounded corners=0.7pt},
}

\begin{axis}[
  width=13.72cm, height=7.3cm, scale only axis,
  xmin=0, xmax=5.4, ymin=-3.4, ymax=0,
  xlabel={$\Rey\omega$}, ylabel={$\Imy\omega$},
  label style={font=\small}, tick label style={font=\scriptsize},
  xtick={0,1,2,3,4,5}, minor xtick={0.5,1.5,2.5,3.5,4.5},
  ytick={0,-1,-2,-3}, minor ytick={-0.5,-1.5,-2.5},
  tick align=outside,
  axis line style={draw=black!55, line width=0.3pt},
  tick style={black!55},
  colormap name=emberdark,
  point meta min=0, point meta max=\atlasMetaMax,
  colorbar,
  colorbar style={
    width=0.22cm, ytick={0,1,2},
    ylabel={$\log_{10}(\text{entries}+1)$},
    ylabel style={font=\scriptsize},
    yticklabel style={font=\scriptsize},
    axis line style={draw=black!55, line width=0.3pt},
    tick style={draw=black!55},
  },
]

% The occupancy field is pixel data and ships as an image.
\addplot graphics [xmin=0, xmax=5.4, ymin=-3.4, ymax=0]
  {data/atlas_density.png};

% gives the colorbar its range without drawing anything
\addplot[scatter, only marks, mark size=0pt, point meta=explicit,
         forget plot, draw opacity=0, fill opacity=0]
  table[x=x, y=y, meta=m] {
    x  y     m
    0  -3.4  0
    0  -3.4  \atlasMetaMax
  };

% ---- overtone ladders, at the median damping of each band -------------
\node[band] at (axis cs:0.07,-0.329) {$N=0$};
\node[band] at (axis cs:0.07,-1.007) {$N=1$};
\node[band] at (axis cs:0.07,-1.726) {$N=2$};
\node[band] at (axis cs:0.07,-2.511) {$N=3$};

% ---- the two windows opened in Figure 2 -------------------------------
% Boxes are the panel windows themselves; each label is tied to its box by a
% short leader so it can sit clear of the occupied cells.
\draw[win] (axis cs:1.425,-0.625) rectangle (axis cs:2.025,-0.025);
\draw[white, line width=0.4pt] (axis cs:2.025,-0.325) -- (axis cs:2.16,-0.325);
\node[winlab, anchor=west] at (axis cs:2.17,-0.325) {Fig.~\ref{fig:pseudo}(a)};

\draw[win] (axis cs:2.459,-1.723) rectangle (axis cs:3.059,-1.123);
\draw[white, line width=0.4pt] (axis cs:3.059,-1.423) -- (axis cs:3.19,-1.423);
\node[winlab, anchor=west] at (axis cs:3.20,-1.423) {Fig.~\ref{fig:pseudo}(b)};
\end{axis}
\end{tikzpicture}

\caption{\textbf{The stored quasinormal atlas.} Occupancy of all
$\atlasN$ stored mode entries, across seven azimuthal sectors and
$\atlasBranches$ continued branches, on a $\atlasGridRe\times\atlasGridIm$ grid
of the complex-frequency plane; the cell value is
$\log_{10}(\text{entries}+1)$, so an empty cell is exactly empty. The overtone
ladders $N=0,1,2,3$ are labelled at their median damping. Boxes mark the two
windows opened in Figure~\ref{fig:pseudo}. Only $\atlasOccupied\%$ of the
plotted cells contain any stored entry: the search is finite, and the figure
shows exactly where it went.}
\label{fig:atlas}
\end{figure}

\begin{table}[t]
\centering
\caption{Finite-atlas census and continuation bookkeeping. Counts refer to
stored computations; they do not constitute a covering certificate for the
continuous parameter domain.}
\label{tab:census}
\begin{tabular}{lc}
\toprule
\thead{Quantity} & \thead{Value} \\
\midrule
  Stored mode entries & 50,558 \\
  Multi-branch parameter points & 4,459 \\
  Independent azimuthal sectors & 7 \\
  Angular branches per sector & 3 \\
  Overtones & $N=0,1,2,3$ \\
  Maximum stored continued-fraction residual & $3.6\times10^{-9}$ \\
  Predictor misses & 118 \\
  Nonconvergences & 33 \\
  Branch-label collapses excluded & 4 \\
\bottomrule
\end{tabular}

\end{table}

\subsection{What the atlas establishes}

Over the stored sample, the extrema recorded in the numerical summary are
\begin{equation}
\dmin=0.5116,\qquad \kappa^{(\text{fixed scaling})}_{\max}=8.73\times10^3 .
\end{equation}
The minimum occurs on the boundary of the sampled region. Relative to the
stated control value $18.3$, the condition indicator is larger by approximately
\begin{equation}
\frac{8.73\times10^3}{18.3}\simeq4.77\times10^2 ,
\end{equation}
not $476$ exactly. This elevation is consistent with a close avoided crossing in
the chosen matrix normalization, but it is not a normalization-independent
distance bound.

Figure~\ref{fig:pseudo} shows what that elevation looks like as a geometric
object. The resolvent landscape at the tightest stored interaction is visibly
more distorted than at the ordinary control, but the level sets close around a
single tracked mode rather than pinching towards a common point, which is the
signature a genuine coalescence would leave.

\begin{figure}[t]
\centering
% Figure 2 -- pseudospectral landscapes.
% Drawn in PGFPlots from data/pseudo_*.dat (the stored fields) and
% data/pseudocont_*.dat (level sets extracted from those same fields).
% Every axis limit, colour limit and marker position below comes from
% data/pseudo_facts.tex, which scripts/tmlr_figdata.py writes next to the
% data, so the frame can never drift away from the field it encloses.
% The two panels share one window size and one square aspect ratio: their
% sharpness is meant to be compared, and that comparison is only honest if
% the plotting geometry is identical.
\begin{tikzpicture}

\pgfplotsset{
  pseudo/.style={
    width=6.45cm, height=6.45cm, scale only axis,
    view={0}{90},
    xlabel={$\Rey\omega$}, ylabel={$\Imy\omega$},
    label style={font=\small}, tick label style={font=\scriptsize},
    axis line style={draw=black!55, line width=0.3pt},
    tick style={black!55}, tick align=outside,
    colormap/viridis,
    colorbar horizontal,
    colorbar style={
      % far enough below the axis to clear the x label, which shares the gap
      height=0.19cm, at={(0,-0.205)}, anchor=north west, width=6.45cm,
      xtick style={draw=black!55},
      xticklabel style={font=\scriptsize},
      xlabel={$\log_{10}\sigma_{\min}/\sigma_{\max}$},
      xlabel style={font=\scriptsize, yshift=2pt},
      axis line style={draw=black!55, line width=0.3pt},
    },
  },
  % Level sets are drawn twice: a soft dark halo first, then the white line
  % on top.  A single white line disappears into the bright end of viridis
  % and a single dark line disappears into the dark end.
  contourhalo/.style={draw=black, line width=1.0pt, draw opacity=0.16,
                      no markers, unbounded coords=jump},
  contourover/.style={draw=white, line width=0.45pt, draw opacity=0.62,
                      no markers, unbounded coords=jump},
  trackedmode/.style={only marks, mark=*, mark size=1.9pt,
                      mark options={fill=white, draw=black, line width=0.6pt}},
  % the partner ring is drawn dark-then-light so it reads on either end of
  % the colour map, the same trick as the level sets
  partnerhalo/.style={only marks, mark=o, mark size=2.4pt,
                      mark options={draw=black, line width=1.5pt,
                                    draw opacity=0.35}},
  partnermode/.style={only marks, mark=o, mark size=2.2pt,
                      mark options={draw=white, line width=0.9pt}},
}
% node options resolve in /tikz/, so this one cannot live in \pgfplotsset
\tikzset{
  kappabox/.style={font=\scriptsize, text=white, inner xsep=2.6pt,
                   inner ysep=1.8pt, fill=black, fill opacity=0.55,
                   text opacity=1, rounded corners=1pt},
}

% ============================== (a) ordinary control ====================
\begin{axis}[pseudo,
  name=ctrl,
  xmin=\psAxmin, xmax=\psAxmax, ymin=\psAymin, ymax=\psAymax,
  xtick={1.5,1.7,1.9}, minor xtick={1.6,1.8,2.0},
  ytick={-0.6,-0.4,-0.2}, minor ytick={-0.5,-0.3,-0.1},
  point meta min=\psAzmin, point meta max=\psAzmax,
  colorbar style={xtick={-8,-7.5,-7,-6.5,-6,-5.5,-5,-4.5,-4,-3.5,-3,-2.5,-2}},
]
\addplot3[surf, shader=interp, draw=none,
          mesh/rows=\psGrid, mesh/cols=\psGrid]
  table[x=re, y=im, z=logsig] {data/pseudo_ordinary_control.dat};
\addplot[contourhalo] table[x=re, y=im]
  {data/pseudocont_ordinary_control.dat};
\addplot[contourover] table[x=re, y=im]
  {data/pseudocont_ordinary_control.dat};
\addplot[trackedmode] table[x=re, y=im]
  {data/pseudo_mark_ordinary_control.dat};
\node[kappabox, anchor=north west] at (rel axis cs:0.035,0.965)
  {$\kappa_T=18.3$};
\end{axis}

% ========================= (b) tightest interaction =====================
\begin{axis}[pseudo,
  name=tight, at={(ctrl.south east)}, xshift=1.95cm, anchor=south west,
  xmin=\psBxmin, xmax=\psBxmax, ymin=\psBymin, ymax=\psBymax,
  xtick={2.5,2.7,2.9}, minor xtick={2.6,2.8,3.0},
  ytick={-1.7,-1.5,-1.3}, minor ytick={-1.6,-1.4,-1.2},
  point meta min=\psBzmin, point meta max=\psBzmax,
  colorbar style={xtick={-8,-7.5,-7,-6.5,-6,-5.5,-5,-4.5,-4,-3.5,-3,-2.5,-2}},
]
\addplot3[surf, shader=interp, draw=none,
          mesh/rows=\psGrid, mesh/cols=\psGrid]
  table[x=re, y=im, z=logsig] {data/pseudo_tightest_interaction.dat};
\addplot[contourhalo] table[x=re, y=im]
  {data/pseudocont_tightest_interaction.dat};
\addplot[contourover] table[x=re, y=im]
  {data/pseudocont_tightest_interaction.dat};
\addplot[trackedmode] table[x=re, y=im]
  {data/pseudo_mark_tightest_interaction.dat};
\addplot[partnerhalo] table[x=re, y=im] {data/pseudo_pair.dat};
\addplot[partnermode] table[x=re, y=im] {data/pseudo_pair.dat};
\node[kappabox, anchor=north west] at (rel axis cs:0.035,0.965)
  {$\kappa_T=8.73\times10^{3}$};
\end{axis}

\node[anchor=south west, font=\bfseries\small, yshift=2pt]
  at (ctrl.north west)  {(a)\; ordinary control};
\node[anchor=south west, font=\bfseries\small, yshift=2pt]
  at (tight.north west) {(b)\; tightest stored interaction};
\end{tikzpicture}

\caption{\textbf{Pseudospectral landscape of the quasinormal pencil.}
$\log_{10}\sigma_{\min}/\sigma_{\max}$ of the Solver~C exterior-complex-scaled
collocation matrix on $\psGrid\times\psGrid$ complex-frequency grids, with
level sets extracted from the same fields. Each panel is centred on its own
tracked mode and both use the same window size and the same square aspect
ratio, so the two landscapes are directly comparable. (a) An ordinary control
at $s=0.24$, $\delta=0.10$, $\mu=0.6$ in sector $(1,1)$, $\ell=2$. (b) The
tightest stored interaction, at the corner $s=0.42$, $\delta=-0.20$, $\mu=1.8$
in sector $(2,0)$, $\ell=4$, whose condition indicator is the atlas maximum
$\kappa_T=8.73\times10^3$. Filled markers are the tracked modes. The open
marker in~(b) is the nearer member of the closest stored pair; its partner lies
below the plotted window at $\Imy\omega=\psPartnerIm$, the two being separated
by $\dmin=\psGap$. Note that the two colour scales differ: the field at~(b)
runs about two decades deeper than at~(a) over a comparable span, so the
landscape there is far more sensitive. Its level sets nonetheless remain
organized around one mode rather than pinching towards a shared point. Both
fields are properties of the finite-dimensional discretization; no continuum
pseudospectral claim is made.}
\label{fig:pseudo}
\end{figure}

The ten closest stored pairs remain separated under the documented local
refinements and show neither one-loop branch exchange nor $t^{1/2}$ splitting.
These observations favor avoided crossings. Candidate-by-candidate
recurrence-free confirmation is not available for every member of this set, so
the classification is kept provisional. The numerical conclusion is therefore:
no exceptional point was detected among the sampled and locally refined
candidates. This statement is restricted to the computed set and does not
exclude EPs throughout the continuous domain in Eq.~\eqref{eq:box}.

\section{Threshold and quasiresonant behavior}

A separate continuation follows eleven branches in six sectors as $\mu$
increases. Along one branch, the damping decreases from $0.357$ at $\mu=0$ to
$8.65\times10^{-10}$ at $\mu=2.968$, while $\Rey\omega\to\mu$. The neighboring
overtone remains at damping of order $10^{-1}$, so the observed behavior is not
a two-mode coalescence in the tracked window.

Exterior complex scaling becomes ineffective in this limit because
Eq.~\eqref{eq:scaling} requires
\begin{equation}
\vartheta_{\min}\to90^\circ \qquad\text{as}\qquad \Rey\Omega\to0 .
\end{equation}
This is a limitation of the two recurrence-free formulations used here, not a
proof that no recurrence-free or hyperboloidal method can treat the threshold.
A hyperboloidal formulation is a promising route into this region
\citep{jaramillo2021,macedo2026}.

No full-atlas superradiant margin can be inferred from the aggregate summary
statistics alone. In particular, Table~\ref{tab:horizons} gives
$\Omega_a=\Omega_b=0.5445$ at $s=0.42$, $\delta=0$, so the $(1,1)$ critical
frequency there is approximately $1.089$, not bounded by $0.51$. Superradiant
statements must be evaluated branch by branch and must distinguish scattering
QNMs from decaying quasibound states.

\section{Physical interpretation and limitations}

The five-dimensional asymptotically flat Myers--Perry black hole is not an
astrophysical model of observed four-dimensional black holes. Its value here is
structural. The sector theorem applies whenever a valid QNM realization
commutes with surviving azimuthal symmetries. In four-dimensional Kerr, for
example, modes of different $m$ lie in distinct invariant sectors; the known
massive-scalar EPs occur within a fixed $m$ sector
\citep{cavalcante2024,cavalcante2026}.

The higher-dimensional instability context must also be stated carefully.
Equal-angular-momentum Myers--Perry black holes were found to show no
gravitational instability in five or seven dimensions in the axisymmetric
analysis of \citet{dias2010}; the instability in that axisymmetric analysis
first appeared in nine dimensions. Other nonaxisymmetric and ultraspinning
instabilities occur in higher-dimensional families and dimensions
\citep{hartnett2013,emparan2014}, but the scalar QNM calculation here does not
establish gravitational stability or instability of the five-dimensional
background.

The principal limitations are:
\begin{enumerate}
\item the theorem is conditional on a specified local holomorphic Fredholm
  realization and excludes only cross-sector mechanisms;
\item the numerical search is finite and the radial recurrence has no certified
  truncation/aliasing enclosure;
\item overtones $N>3$, higher angular branches, $s>0.42$, $\mu>1.8$ in the
  atlas, and the superradiant bound-state regime were not searched;
\item the threshold region is not recurrence-free verified by the formulations
  used here.
\end{enumerate}

\section{Discussion and conclusions}

The central conclusion is structural. Equal-magnitude-spin symmetry produces
exact frequency degeneracies, but the surviving azimuthal torus symmetry
simultaneously prevents states in different $(m_1,m_2)$ sectors from forming a
shared defective block. Exceptional behavior therefore remains possible only
through collisions of distinct spectral branches inside one fixed azimuthal
sector.

No such exceptional collision was detected in the finite atlas. The smallest
stored gap is $0.5116$, and the closest candidates show neither branch exchange
nor square-root splitting under the documented local tests. Because the
calculation does not include a certified covering or radial-truncation
enclosure, this result is a finite-atlas non-detection, not a theorem excluding
exceptional points throughout the continuous search box.

Two methodological lessons are broadly useful. First, derivatives and Taylor
coefficients of an arbitrarily normalized meromorphic spectral condition are
not absolute proximity measures. Second, a small gap between numerical labels is
evidence about branch tracking before it is evidence about physics;
continuation histories, independent formulations, monodromy, and local Puiseux
fits are necessary safeguards.

\subsubsection*{Data availability}

The numerical continuation atlas and solver materials supporting the
finite-atlas claims are available from the authors upon reasonable request
through the anonymized review process. They are not publicly archived at the
time of submission. The anonymized Supplementary Material contains the solver
inventory, atlas census, verification gate, symmetry checks, and threshold
summary.

\bibliographystyle{tmlr}
\bibliography{refs}

\appendix

\section{Checks on separation and singular structure}

The determinant of Eq.~\eqref{eq:metric} satisfies
\begin{equation}
\sqrt{|g|}=r\,\Sigma\sin\theta\cos\theta .
\end{equation}
Substitution of Eq.~\eqref{eq:ansatz} into
\begin{equation}
\partial_\nu\bigl(\sqrt{|g|}\,g^{\nu\rho}\partial_\rho\Phi\bigr)
=\mu^2\sqrt{|g|}\,\Phi
\end{equation}
separates into Eqs.~\eqref{eq:angular} and \eqref{eq:radial}. In $z=r^2$, the
potential has apparent singularities at $z=0,z_+,z_-$. When $ab\neq0$, the
apparent $1/z$ terms cancel because $W(0)=abG$ and $P(0)=a^2b^2$. The physical
regular singular points are $z_\pm$, and infinity is an irregular singular point
of rank one.

At $z=z_+$, factorization gives
\begin{equation}
W(z_+)=(z_++a^2)(z_++b^2)(\omega-\omega_c),
\end{equation}
leading to Eq.~\eqref{eq:horizonbc}. At infinity, division of
Eq.~\eqref{eq:radial} by its leading radial coefficient gives the asymptotic
equation preceding Eq.~\eqref{eq:infinitybc}. The absence of an $O(r^{-1})$ term
removes the four-dimensional Coulomb logarithm.

\section{Keldysh pole order}

Let $T(\omega)$ be holomorphic and Fredholm of index zero in a neighborhood of
an isolated characteristic value $\omega_*$. For a simple eigenvalue, with right
and left null vectors normalized by $\langle v,T'(\omega_*)u\rangle=1$,
\begin{equation}
T(\omega)^{-1}=\frac{u\otimes v}{\omega-\omega_*}+O(1).
\end{equation}
For maximal chain length $q$, the local Keldysh expansion has
\begin{equation}
T(\omega)^{-1}=\sum_{k=1}^{q}\frac{S_k}{(\omega-\omega_*)^k}+O(1),
\qquad S_q\neq0 .
\end{equation}
Thus the pole order equals the maximal partial multiplicity. Inverse Laplace
transformation of a pole of order $k$ produces
$t^{k-1}e^{-i\omega_*t}/(k-1)!$, giving Eq.~\eqref{eq:keldysh}. Combined with
the sector lemma, this proves the stated relation between cross-sector
coincidences and ringdown pole order.

\end{document}
